%----------------------------------------------------------------------------
\chapter{\bevezetes}
\label{chap:Introduction}
%----------------------------------------------------------------------------

%----------------------------------------------------------------------------
\section{Motivation}
\label{sec:Motivation}
%----------------------------------------------------------------------------
\begin{itemize}
    \item Real-time simulation is a baseline expectation in games, VR, and AR.
    \item The simulator must deliver believable dynamics within a tight frame budget. A 60~fps framerate allows $\sim$16~ms for all physics, rendering, and gameplay logic, with physics getting only a fraction.
    \item I focus on deformable bodies (springs, chains and cloth) rather than rigid bodies because:
    \begin{itemize}
        \item Deformable simulation is the harder case (many more degrees of freedom and constraints per frame), stressing both solvers.
        \item Even with vertices and constraints in the ten thousands, all parts of the simulated system are still visible.
        \item Solver behaviour can be easily observed, failures can be spotted.
        \item Springs are one of the most common primitives in simulations, and an easily measureable one. Cloth is also really common, and a crucial stepping stone to soft body simulation.
    \end{itemize}
\end{itemize}

Designing such a (real-time) simulator is hard for three recurring reasons:
\begin{itemize}
    \item \emph{Stiffness.} Cloth and soft tissue have very stiff constraints (almost-inextensible fabric, near-rigid bones); stiff systems make explicit integrators unstable unless time steps become impractically small.
    \item \emph{Stability.} Even implicit integrators can blow up under poorly conditioned constraints, large mass ratios between connected particles, or contact-induced velocity discontinuities.
    \item \emph{Strict time constraints.} Methods that converge well offline (Newton's method, projective dynamics) are too expensive per frame; cheap real-time methods often trade convergence quality for predictable cost.
\end{itemize}

\paragraph{The accuracy/efficiency trade-off.}
\begin{itemize}
    \item Considering the above, every physics-based simulator sits on a line between \emph{physical accuracy} (merging stiffness and stability) and \emph{computational efficiency}. Trying to optimize efficiency will result in less realism.
    \item \autoref{tab:accuracy-efficiency} sketches this spectrum with the most known solvers for deformable body simulation, ordered from most accurate to most efficient.
    \item Mass-Spring System (MSS) is fast. It represents deformable materials using a network of masses connected by springs and uses explicit Euler integration, the most intuitive and straight-forward approach. It's technically more physically accurate than PBD or XPBD, but it is highly unstable for stiff and complex systems, sending it to the bottom of the list. \emph{TODO: cite something that shows this}
    \item On the other end of the spectrum, the Finite Element Method (FEM) excels in high physical accuracy by trating objects as a continuous material discretized into volumetric finite elements (such as tetrahedra). It works with a Newton solver (using implicit Euler integration). However, its really high computational cost makes it impractical for real-time applications. \cite{nikishkov2004introfem}
\end{itemize}

\begin{table}[h]
\centering
\begin{tabular}{l l l}
\toprule
\textbf{Method} & \textbf{Accuracy} & \textbf{Cost} \\
\midrule
FEM     & reference            & most expensive \\
VBD     & moderate--high       & moderate       \\
XPBD    & moderate             & cheap          \\
PBD     & low (stiffness-dep.) & cheap          \\
MSS     & poor (unstable)      & cheapest       \\
\bottomrule
\end{tabular}
\caption{Deformable-simulation methods ordered from most accurate to most
efficient. XPBD and VBD the real-time domain.}
\label{tab:accuracy-efficiency}
\end{table}

XPBD and VBD both sit near the efficient end, within the real-time cathegory.
\begin{itemize}
    \item XPBD~\cite{10.1145/2994258.2994272} has been a workhorse for years.
    \item VBD~\cite{Chen2024} is a recent, promising addition not yet throughly evaluated head-to-head against XPBD.
\end{itemize}

%----------------------------------------------------------------------------
\section{Contribution}
\label{sec:Contribution}
%----------------------------------------------------------------------------
\begin{itemize}
    \item The central contribution is, to the best of my knowledge, the \emph{first faithful XPBD-vs-VBD comparison}:
    \begin{itemize}
        \item both algorithms implemented in the same Unity-based codebase,
        \item applied to the same mass-spring material model,
        \item run under a matched compute budget (iteration and substep counts), not merely matched material parameters,
        \item evaluated not only qualitatively but also against a high-iteration Newton solver serving as a numerical ground-truth reference, which lets me report \emph{solution error} rather than only relative performance.
    \end{itemize}
    \item Crucially, the comparison uses XPBD --- the constraint-stiffness-independent variant any practitioner would deploy today --- and the canonical VBD formulation, not plain PBD or a non-canonical variant.
    \item This sharpens the closest prior comparison, Ma et~al.~\cite{app142311072}, which contrasts plain PBD with a non-canonical VBD implementation and reports only frame-rate scaling; \autoref{sec:Ma2024} sets out exactly how my setup differs.
\end{itemize}

%----------------------------------------------------------------------------
\section{The primal/dual lens}
\label{sec:PrimalDualLens}
%----------------------------------------------------------------------------
\begin{itemize}
    \item The two solvers sit on opposite sides of the \emph{primal/dual} duality identified by Macklin et~al.~\cite{10.1111/cgf.14104}: XPBD is the dual solver (it solves for Lagrange multipliers on constraints), VBD the primal one (it solves directly for vertex positions).
    \item This duality is the organising lens of the thesis: it anticipates complementary weak spots --- the dual solver where information must propagate along chains of coupled constraints, the primal solver where a wide spread of stiffnesses meets at one vertex --- and the stress tests in \autoref{chap:Experiments} are built around exactly these scenarios. \autoref{sec:PrimalDual} develops the distinction and its consequences in full.
\end{itemize}

%----------------------------------------------------------------------------
\section{Evaluation axes}
\label{sec:EvalAxes}
%----------------------------------------------------------------------------
I evaluate both solvers along four axes that together capture the trade-offs a practitioner cares about:
\begin{description}
    \item[Efficiency.] Computational cost per frame as a function of mesh size, substep/iteration count.
    \item[Robustness.] Behaviour under stress: extreme mass ratios, extreme stiffness ratios, and how they result in stability issues.
    \item[Physical realism.] Energy behaviour, damping behaviour, and how close each method is from a ground truth Newton reference.
    \item[Flexibility.] Checking independency of material properties from iteration counts, adding new constraint types, and considering the domain coverage of the models.
\end{description}

%----------------------------------------------------------------------------
\section{Thesis structure}
\label{sec:Structure}
%----------------------------------------------------------------------------
\begin{itemize}
    \item \autoref{chap:Background}: theoretical foundation --- variational form of implicit Euler, mass-spring energies, the primal/dual distinction.
    \item \autoref{chap:RelatedWork}: reviews PBD, XPBD, VBD literature and positions this work against Ma et~al.~\cite{app142311072}.
    \item \autoref{chap:Algorithms}: states the two algorithms in canonical form.
    \item \autoref{chap:Methodology}: the Unity implementation and the design decisions that keep the comparison fair.
    \item \autoref{chap:Experiments}: results along the four axes.
    \item \autoref{chap:Discussion}: evaluating results, limitations.
    \item \autoref{chap:Conclusion}: closing with future-work directions.
\end{itemize}
